\clearpage\thispagestyle{empty}
\pdfbookmark[0]{Quellen zur Rezension der Connoissance des tems für 1810}{connoissance1810-sources}
\begin{center}\textbf{Quellen zur Rezension der Connoissance des tems für 1810}\end{center}
\noindent Diese KI-übersetzte Fassung in modernem Deutsch folgt Carl Friedrich Gauß, \href{https://gdz.sub.uni-goettingen.de/id/PPN236007467}{\emph{Werke}, Band~VI} (Göttingen, 1874), S.~515–517. Die Rezension erschien zuerst ohne Verfassernamen unter \emph{Paris} in den \emph{Göttingischen gelehrten Anzeigen}, Stück~95, 17. Juni 1809, S.~939–943. Die Niedersächsische Staats- und Universitätsbibliothek Göttingen stellt den \href{https://images.p.sub.uni-goettingen.de/iiif/image/gdz:PPN319721507_1809:00000941/full/1650,/0/default.jpg}{Anfang}, den \href{https://images.p.sub.uni-goettingen.de/iiif/image/gdz:PPN319721507_1809:00000943/full/1650,/0/default.jpg}{Abschnitt zur Gradmessung} und den \href{https://images.p.sub.uni-goettingen.de/iiif/image/gdz:PPN319721507_1809:00000945/full/1650,/0/default.jpg}{Schluss} bereit. Der Erstdruck gibt für die zweite Abplattung $1/317$ und für den zweiten Ortsnamen \emph{Trivandeporum} an. Die Werkausgabe hat $1/817$ und \emph{Trivandepcrum}. Diese Lesarten bleiben hier erhalten. Im zweiten Breitengradwert steht im Erstdruck ein Komma und in der Werkausgabe ein Punkt. Am Anfang steht im Erstdruck \emph{geographischen}, in der Werkausgabe \emph{geopraphischen}.
\par\smallskip\noindent Der französische Titel bedeutet \emph{Kenntnis der Zeiten oder der Himmelsbewegungen zum Gebrauch der Astronomen und Seefahrer für das Jahr 1810}. Das Bureau des longitudes gab den \href{https://cdt.imcce.fr/items/show/159}{Jahrgang für 1810} heraus. Gedruckt wurde er 1808 in der Kaiserlichen Druckerei in Paris. Zu den verwandten Arbeiten in diesem Band gehören Gauß’ \hyperlink{native-128}{\emph{Allgemeine Tafeln für Aberration und Nutation}} und seine \hyperlink{native-517-lower}{Rezension des astronomischen Teils der \emph{Reise von Humboldt und Bonpland}}. Übersetzung und Satz mit KI erstellt.
\clearpage
